This section reviews the prior work most relevant to this study, spanning
SOH prediction methods, large-format and LFP cells, transfer learning,
manufacturing variability, and uncertainty quantification.

\subsection{SOH Prediction Methods}

Data-driven SOH prediction has evolved from simple regression models to
sophisticated deep-learning architectures.
Recurrent neural networks, particularly LSTM variants, have emerged as
the dominant approach due to their ability to model long-range temporal
dependencies in discharge capacity sequences~\cite{zhang2018long}.
Convolutional architectures form a second major family: deep CNNs with
ensemble and transfer learning have been applied to capacity
estimation~\cite{shen2022deep,li2020lithium}, building on early
demonstrations that deep networks can estimate battery state directly from
operational measurements~\cite{chemali2018state}.

Attention mechanisms have been incorporated to allow models to weight
different timesteps in the input window~\cite{bahdanau2015neural}.
Qu et al.\ proposed PA-LSTM, which augments an LSTM encoder with an
attention layer and uses particle swarm optimisation (PSO) for hyperparameter
tuning and model pretraining, with incremental learning for online SOH
monitoring; PA-LSTM outperformed vanilla LSTM, recurrent neural network (RNN), and relevance vector machine (RVM) baselines on
NASA cells~\cite{qu2019neural}.
Subsequent work replaced PSO with other metaheuristics for LSTM hyperparameter
optimisation: genetic algorithm (GA-LSTM)~\cite{ponnambalam2024battery}
and differential evolution~\cite{storn1997differential}
(DE-LSTM)~\cite{ponnambalam2026delstm} both
achieved significant RMSE reductions on the same NASA cells (B05, B06, B18),
retained from PA-LSTM to enable direct comparison.

A common thread across these methods is the use of \textit{incremental
(walk-forward) learning}: the model is trained on the first $p$\% of a cell's
SOH trajectory and evaluated on the remainder, simulating online deployment.
This protocol is adopted throughout the present work.
The proposed hybrid uses the same walk-forward protocol but introduces a
fundamentally different use of early-cycle data: rather than fitting the
observed sequence to continue it, a convolutional encoder summarises the
\textit{shape} of the first 30 cycles into a cell-identity embedding that
initialises the LSTM hidden state before fine tuning begins.
This places each cell in the pretrained population's manufacturing
distribution, discriminating between cell types even when early trajectory
magnitudes appear similar to a sequence-fitting model.

\subsection{Large-Format and LFP Batteries}

Most published SOH prediction studies use 18650-format cylindrical cells
(LCO, NMC, NCA, or LFP chemistry, 1 to 3\,Ah).
Large-format prismatic and pouch cells are the dominant format in stationary
storage and electric heavy vehicles~\cite{barai2015comparison}, yet they
are substantially underrepresented in the prediction literature.
Key differences include: non-uniform current distribution across the
electrode area~\cite{forgez2010thermal}; larger absolute capacity changes
per SOH unit (a 1\% SOH change in a 50\,Ah cell = 0.5\,Ah vs.\ 0.02\,Ah
in a 2\,Ah cell); and different thermal management requirements.

LiFePO$_4$ (LFP) cells are preferred for stationary storage due to their
flat discharge plateau, long cycle life, and superior thermal
stability~\cite{barai2015comparison}.
The flat voltage plateau of LFP makes capacity based SOH estimation more
robust than voltage based methods, motivating the discharge-capacity
incremental learning approach used in this work.

\subsection{Transfer Learning for Battery SOH}

Transfer learning (pretraining a model on a source dataset and adapting
it to a target domain) has attracted growing interest for battery
prognostics~\cite{severson2019data}.
Approaches include fine tuning pretrained networks, domain-adversarial
training, and meta-learning~\cite{finn2017model}.
Yu et al.~\cite{yu2026crossdomain} addressed cross-domain SOH estimation across
heterogeneous small-cell datasets (NASA and CALCE), fusing the most similar
source domains through a similarity network to maintain accuracy when
target-cell data are scarce.
Ye and Yu~\cite{ye2021state} applied domain-adversarial transfer learning to
align feature distributions between source and target cells for SOH estimation.
Closest to the present work, Ma et al.~\cite{ma2022real} achieve personalised,
real-time prediction by fine tuning a pretrained model on each cell's most
recent cycles.
The fingerprint encoder proposed here uses early-life data differently: it
treats the first 30 cycles as a static shape that sets the LSTM initial state
before fine tuning, rather than fine tuning on those cycles directly
(Section~\ref{sec:mechanism}).

A key practical constraint is chemistry: LFP cells have a fundamentally
different degradation mechanism (lithium plating, phase transition stress)
compared to layered-oxide cells such as LiCoO$_2$ (LCO) and NMC (transition metal dissolution, solid electrolyte interphase (SEI) growth)~\cite{birkl2017degradation}.
Whether a model pretrained on one chemistry can transfer to another
is an open question addressed in this work through explicit comparison
of chemistry matched (LFP$\rightarrow$LFP) and chemistry mismatched
(LFP$\rightarrow$LCO) transfer.

\subsection{Manufacturing Variability}

Cell-to-cell variability arising from manufacturing tolerances is
well-documented in the electrochemical literature~\cite{birkl2017degradation}
but rarely quantified in prediction studies.
Dubarry et al.~\cite{dubarry2018synthesize} showed that even cells from
the same batch exhibit measurable differences in electrode loading and
electrolyte fill that propagate to divergent long-term trajectories.
Severson et al.~\cite{severson2019data} observed that early-cycle
statistics (first 100 cycles) correlate with cycle life, motivating
the fingerprint approach adopted here.

To our knowledge, no prior work has simultaneously quantified
manufacturing variability across 10 identical large-format LFP cells
and studied how cross-cell pooling interacts with this variability,
the \textit{two-sided failure mode} we characterise in Section~\ref{sec:variability}.

\subsection{Uncertainty Quantification}

Point-estimate SOH predictions are insufficient for BMS deployment, where
decisions under uncertainty require calibrated prediction intervals.
Bayesian neural networks~\cite{blundell2015weight} and Monte Carlo
Dropout~\cite{gal2016dropout} have been applied to battery SOH uncertainty
estimation~\cite{ke2024bayesian}, but calibration is rarely verified against
held-out data, and never across a chemistry mismatch.
Conformal prediction provides distribution free coverage
guarantees~\cite{angelopoulos2021gentle}, an attractive property for
safety-critical battery management.
In this work we use MC-Dropout and quantify its empirical coverage, revealing
a chemistry-alignment effect in which calibration holds only when the
pretraining and deployment chemistries match.
